\chapter{Prediction markets and scoring rules}\label{ch:pm}

\noindent A prediction market is the simplest derivative there is: a share that pays \$1 if an
event happens and \$0 otherwise. Its price is a probability estimate with money behind it. For
a researcher the venue is less important than what it produces---a large, public, continuously
updated panel of probabilistic forecasts you can score, calibrate and cross-check. The canonical
venue for this book is Polymarket.

\section{The contract}\label{sec:pmcontract}

A YES share costs $p\in(0,1)$ and resolves to $1$ or $0$. A NO share is the complement: it costs
$1-p$ and resolves to the opposite side. Buying YES at $p$ is a bet with
$\mathrm{EV}=q-p$ per share where $q$ is the event's true probability, and variance
$q(1-q)$---binary outcomes are among the highest-variance instruments that exist per unit of
expected value.

\begin{figure}[htbp]\centering
\begin{tikzpicture}[font=\scriptsize]
  \node[draw=slate,fill=white,rounded corners=1pt,align=center,minimum width=2.6cm] (buy) at (0,0)
    {pay $p$ now\\ e.g.\ $0.62$};
  \node[draw=good,fill=goodbg,rounded corners=1pt,align=center,minimum width=2.6cm] (yes) at (4.4,0.9)
    {event happens\\ receive $1$};
  \node[draw=amber,fill=amberbg,rounded corners=1pt,align=center,minimum width=2.6cm] (no) at (4.4,-0.9)
    {event fails\\ receive $0$};
  \draw[->,ink,thick] (buy) -- (yes) node[midway,above,slate]{prob $q$};
  \draw[->,ink,thick] (buy) -- (no) node[midway,below,slate]{prob $1-q$};
  \node[align=center,text=slate] at (2.2,-1.9)
    {$\mathrm{EV}=q-p$ per share;\quad $\sigma=\sqrt{q(1-q)}$};
\end{tikzpicture}

\caption{Anatomy of a YES share.}
\label{fig:contract}
\end{figure}

\begin{principle}[price is a probability]
In a frictionless prediction market the YES price equals the marginal trader's subjective
probability. Real venues have fees, latency and incomplete books, so read $p$ as
\emph{consensus-minus-frictions}: the best single public number for ``how likely is this,''
not a guaranteed fair price.
\end{principle}

\begin{example}\label{ex:pmev}
\textbf{Expected value.} A market prices YES at $0.62$. You assess $q=0.70$. What is the edge
per share, the variance, and how many shares must you buy before the bet's standard deviation
is below its expected profit?
\begin{solution}
$\mathrm{EV}=0.70-0.62=+0.08$ per share; $\sigma=\sqrt{0.7\cdot0.3}=0.458$ per share. For $n$
independent copies the ratio is $n\cdot0.08/(\sqrt{n}\cdot0.458)=\sqrt{n}\cdot0.175>1$ at
$n>33$ shares. A single binary bet is essentially a coin flip with a weighted coin; edge is a
\emph{property of repetition}, not of one trade.
\end{solution}
\end{example}

\section{Reading a live book}\label{sec:pmlive}

Polymarket's public Gamma API returns market metadata and live outcome prices, no
authentication required. The 100 most-traded active markets are pinned at
\texttt{data/pm\_markets.json}:

\begin{lstlisting}[language=Python, numbers=none]
d=json.load(open("data/pm_markets.json"))
pr=[float(x) for x in json.loads(m["outcomePrices"])]
\end{lstlisting}
\begin{lstlisting}[numbers=none]
Xi Jinping out before 2027? | ["0.0285", "0.9715"]
Will Gavin Newsom win the 2028 Democratic pres. nomination | ["0.1545", ...]
\end{lstlisting}

YES$+$NO sums to $1$ within a few cents everywhere in the snapshot
(\texttt{arb\_check.py}: ``100 binary markets checked, YES+NO deviating $>$3c from 1: 0'').
That is the arbitrage channel doing its job: any persistent gap is a machine-readable coupon.

\begin{figure}[htbp]\centering
\begin{tikzpicture}
\begin{axis}[width=0.86\linewidth,height=0.32\linewidth,
    xlabel={YES price},ylabel={markets},
    xmin=0,xmax=1,ymin=0,ymax=24,
    tick style={draw=none},axis line style={draw=slate},
    xtick={0,0.25,0.5,0.75,1},font=\scriptsize,ybar,bar width=10pt]
  \addplot[fill=accentsoft,draw=accent] file{figures/data/H_PM.tex};
\end{axis}
\end{tikzpicture}

\caption{YES-price histogram over the 100 most-traded Polymarket markets (snapshot
2026-10-11): a U-shape---most action concentrates at the near-certain and near-impossible
tails. \texttt{data/pm\_markets.json}.}
\label{fig:pmhist}
\end{figure}

\section{Scoring rules}\label{sec:scoring}

To judge a probability forecaster you need a score. A scoring rule is \emph{proper} if the
expected score is optimised by reporting your true belief. The two canonical ones:

\begin{principle}[proper scores]
\textbf{Brier:} $S(p,y)=(p-y)^2$ --- bounded in $[0,1]$, tolerant of wild reports.
\textbf{Log score:} $S(p,y)=\log(p)$ if $y{=}1$ else $\log(1-p)$ --- strictly proper,
\emph{unbounded}: reporting $p=0$ on an event that happens is penalised by $-\infty$.
Properness is what makes an honest forecaster indifferent between reporting and predicting.
\end{principle}

\begin{example}\label{ex:proper}
\textbf{Truth is optimal.} \texttt{scoring.py} computes expected scores for reports $p$ when
the truth is $q=0.6$:
\begin{lstlisting}[numbers=none]
q=0.6 report 0.3: E[Brier]=0.3300  E[log]=-0.8651
q=0.6 report 0.5: E[Brier]=0.2500  E[log]=-0.6931
q=0.6 report 0.6: E[Brier]=0.2400  E[log]=-0.6730
q=0.6 report 0.7: E[Brier]=0.2500  E[log]=-0.6956
q=0.6 report 0.9: E[Brier]=0.3300  E[log]=-0.9843
proper: best report = truth 0.6 for both rules
\end{lstlisting}
\begin{solution}
Both rules bottom out at $p=q$: Brier reaches its minimum $0.24$ and log score its maximum
$-0.673$ exactly at the truthful report, and both degrade symmetrically as the report moves
away. That is propriety in one table: \emph{honesty is a dominant strategy}. (The log rule's
asymmetry $-0.6931$ vs $-0.6956$ is small here but explodes for extreme reports---see the
exercise on log penalty.)
\end{solution}
\end{example}

\begin{figure}[htbp]\centering
\begin{tikzpicture}
\begin{axis}[width=0.68\linewidth,height=0.36\linewidth,
    xlabel={reported $p$ (truth $q=0.6$)},ylabel={expected Brier penalty},
    xmin=0,xmax=1,ymin=0.22,ymax=0.36,
    tick style={draw=none},axis line style={draw=slate},font=\scriptsize]
  \addplot[accent,very thick] file{figures/data/S_BRIER.tex};
  \addplot[slate,dashed] coordinates {(0.6,0.22)(0.6,0.36)};
  \node[font=\scriptsize,text=slate] at (axis cs:0.6,0.35) {minimum at $p=q$};
\end{axis}
\end{tikzpicture}

\caption{Expected Brier penalty vs reported probability when $q=0.6$: minimised at the truth.}
\label{fig:score}
\end{figure}

\section{Calibration}\label{sec:calib}

Scores rank forecasters; \emph{calibration} diagnoses them. A forecaster is calibrated if,
among all events it assigned probability $p$, the empirical frequency is $p$. The standard
tool is the reliability diagram: bin forecasts, plot mean forecast vs empirical rate, and
compare with the diagonal. The summary statistic is expected calibration error,
$\mathrm{ECE}=\sum_b \frac{n_b}{n}\lvert \bar f_b - \bar y_b\rvert$.

\begin{example}\label{ex:calib}
\textbf{Overconfidence detector.} \texttt{calibration.py} builds an overconfident forecaster
(real outcomes compressed toward 0.5 relative to the issued forecast) and bins 4000
predictions:
\begin{lstlisting}[numbers=none]
bin   mean_forecast  empirical   n
0.0-0.1  0.043          0.239      134
0.4-0.5  0.452          0.469      676
0.7-0.8  0.746          0.642      383
0.9-1.0  0.954          0.857      133
ECE=0.0747  (overconfident: forecasts more extreme than reality)
calibrated twin ECE=0.0230
\end{lstlisting}
\begin{solution}
Events forecast at 4\% actually happen 24\% of the time; events forecast at 95\% happen only
86\%---both tails bend \emph{toward the middle}, the fingerprint of overconfidence. The
calibrated twin (forecaster $=$ true generator) scores ECE $0.023$, mostly binning noise.
ECE is not a score---it cannot rank a bad-but-calibrated forecaster against a good one---but
as a diagnostic it separates \emph{miscalibration} from \emph{lack of skill}.
\end{solution}
\end{example}

\begin{figure}[htbp]\centering
\begin{tikzpicture}
\begin{axis}[width=0.62\linewidth,height=0.62\linewidth,
    xlabel={forecast probability},ylabel={empirical frequency},
    xmin=0,xmax=1,ymin=0,ymax=1,
    tick style={draw=none},axis line style={draw=slate},font=\scriptsize]
  \addplot[slate,dashed,domain=0:1] {x};
  \addplot[accent,very thick,mark=*,only marks,mark size=2.2pt] file{figures/data/R_REL.tex};
  \node[font=\scriptsize,text=slate,align=left] at (axis cs:0.16,0.82)
    {overconfident:\\ forecasts more\\ extreme than\\ reality};
\end{axis}
\end{tikzpicture}

\caption{Reliability diagram of the overconfident forecaster: the curve bows below the
diagonal at high forecasts and above it at low ones. \texttt{calibration.py}.}
\label{fig:reliability}
\end{figure}

\section{Consistency and the arbitrage channel}\label{sec:pmarb}

Mutually exclusive outcomes must price to $1$ (plus venue vig), or a Dutch book exists:
buy all outcomes when $\sum p < 1$ for a guaranteed payout of $1$, or sell all when
$\sum p>1$. In practice the gap rarely exceeds a few cents on liquid markets---our snapshot
found zero violations beyond 3c---but the \emph{threat} of the trade is what keeps prices
coherent, which is why summed prices are a live health check.

\begin{figure}[htbp]\centering
\begin{tikzpicture}[font=\scriptsize,every node/.style={align=center}]
  \node[draw=slate,fill=white,rounded corners=1pt,minimum width=3.2cm] (m) at (0,0)
    {three outcomes\\ $p_A{+}p_B{+}p_C=0.97$};
  \node[draw=good,fill=goodbg,rounded corners=1pt,minimum width=3.2cm] (a) at (4.2,0.8)
    {buy all three\\ pay $0.97$, get $1$};
  \node[draw=amber,fill=amberbg,rounded corners=1pt,minimum width=3.2cm] (b) at (4.2,-0.8)
    {sum $>1$: sell all\\ collect $1$, pay $1$};
  \draw[->,ink] (m) -- (a);
  \draw[->,ink] (m) -- (b);
  \node[text=slate] at (2.4,-1.7) {mutually exclusive outcomes must sum to $1$ (minus vig)};
\end{tikzpicture}

\caption{The arbitrage channel for multi-outcome markets.}
\label{fig:arb}
\end{figure}

\begin{example}\label{ex:tails}
\textbf{Reading the tails.} In our real snapshot, 21 of 100 markets priced the lead outcome
below 5c---most of them sports spreads, e.g.\ ``Spread: Viking FK $(-1.5)$'' at 0.9c.
\begin{solution}
Tail prices are where the friction model matters most: a 1c share pays 100:1, so a tiny
mispricing is a huge \emph{relative} edge---and exactly where books are thinnest and
resolutions most ambiguous. The U-shaped histogram (fig.\ \ref{fig:pmhist}) says the median
market is already near resolution; the information density lives in the 5--95c band where
probabilities are actually contested.
\end{solution}
\end{example}

\begin{example}\label{ex:respath}
\textbf{Price paths converge to truth or die trying.} \texttt{resolve\_sim.py} runs a running
estimate over 200 noisy observations of a $q=0.6$ event (fig.\ \ref{fig:path}):
\begin{lstlisting}[numbers=none]
seed1: empirical freq drifts to 0.59 -> resolves 1
seed2: empirical freq drifts to 0.57 -> resolves 1
\end{lstlisting}
A market's price walks the same way: information arrives, the estimate concentrates, and the
price either ratchets toward 1 or collapses to 0. The interesting trades are mid-path---before
convergence is priced in.
\end{example}

\begin{figure}[htbp]\centering
\begin{tikzpicture}
\begin{axis}[width=0.86\linewidth,height=0.32\linewidth,
    xlabel={observation},ylabel={running estimate},
    xmin=0,xmax=200,ymin=0.3,ymax=0.9,
    tick style={draw=none},axis line style={draw=slate},
    xtick={0,50,100,150,200},font=\scriptsize]
  \addplot[accent,thick] file{figures/data/P_PATH.tex};
  \addplot[slate,dashed] coordinates {(0,0.6)(200,0.6)};
  \node[font=\scriptsize,text=slate] at (axis cs:150,0.66) {true $q=0.6$};
\end{axis}
\end{tikzpicture}

\caption{A running empirical frequency over 200 observations of a $q=0.6$ event---the shape
every market price follows toward resolution. \texttt{resolve\_sim.py}.}
\label{fig:path}
\end{figure}

\begin{checkpoint}[Score a forecaster]
\begin{itemize}
\item Compute Brier and log score for reports $(0.8,0.6,0.3)$ on events $(1,0,1)$; who won?
\item From \texttt{data/pm\_markets.json}, find any market whose YES$+$NO deviates $>1$c and
explain what a bot does about it.
\item Why does ECE fail as a \emph{ranking} metric? (A constant 0.5 forecaster can be perfectly
calibrated on symmetric data yet have no skill.)
\end{itemize}
\end{checkpoint}

\section{Recap: the pipeline}\label{sec:pmrecap}

Fetch prices, treat each as a probability, watch consistency for arbitrage, score predictions
with a proper rule, and diagnose with calibration. Everything in later chapters---backtesting,
signal evaluation, sizing---reuses this loop: markets emit probabilities; proper scores tell
you whether the emitter was any good.

\begin{example}[The guarantee that costs \$1.03]\label{ex:lock}
\textbf{Strategy.} YES$+$NO $=\$1$ at settlement, so buying one of
each is a riskless \$1. Price that portfolio on every live book in
the Polymarket snapshot, using the ask side only (you must lift,
not rest).

\textbf{Solution.} \texttt{book\_width.py} prints:
\begin{verbatim}
markets with a live two-sided book: 40/100
bid-ask width: median 3.0c mean 12.5c
tightest 0.1c widest 94.0c
YES+NO bought at ask: median $1.030  <$0.99: 0  >$1.01: 30
cheapest lock cost $1.001 -> guaranteed $-0.001 per $1
\end{verbatim}
Only 40 of 100 markets even quote both sides. On those, the
riskless bundle costs a median $\$1.03$ for a guaranteed \$1 ---
the spread \emph{is} the venue's edge, and no market lets you lock
a profit. Even the tightest book leaves you $-0.1$c down.

\textbf{Significance.} Mid-price arithmetic says YES$+$NO $=1$
everywhere; the \emph{actionable} price says it costs $\ge1.001$.
Every PM ``arbitrage'' you scan in outcomes data is the spread
wearing a disguise. This is the same lesson as the LOB chapter,
one market microstructure up.
\end{example}

\begin{example}[How much does the market know?]\label{ex:conviction}
\textbf{Strategy.} Treat each market's YES-price as a forecast and
ask the Brier score what it \emph{expects} to score if the market
were perfectly calibrated. That converts 100 price ticks into an
information budget.

\textbf{Solution.} \texttt{scoring\_stakes.py} prints:
\begin{verbatim}
conviction |2p-1| over 100 binary markets
mean 0.553   >0.9: 25   <0.1: 20
E[Brier] if market calibrated = 0.1386 vs 50c baseline 0.25
info gain over coin-flip = 45%
\end{verbatim}
A quarter of the board prices near-certainty ($|2p-1|>0.9$) and a
fifth sits near $50/50$ --- the market is either very sure or
barely guessing, rarely in between. If it were calibrated, its
expected Brier ($0.139$) beats the coin-flip baseline by $45\%$.

\textbf{Significance.} Calibration is a property of the
\emph{population of prices}, not one market: $E[p(1-p)]$ prices the
information content of the whole board in one number. It also
quantifies the stakes --- at $45\%$ better than guessing, a
forecaster who is merely ``as calibrated as the board'' already
carries real edge.
\end{example}

\section{The automated alternative: LMSR}\label{sec:lmsr}
Every market so far was an order book. Prediction markets grew up
on a different mechanism: the \emph{logarithmic market scoring
rule}. The operator keeps a cost function
$C(q_Y,q_N)=b\log\!\big(e^{q_Y/b}+e^{q_N/b}\big)$ where $q_i$ are
shares sold and $b$ is a liquidity parameter; the instantaneous
YES-price is the share of ``mass'' on YES,
$p=\frac{e^{q_Y/b}}{e^{q_Y/b}+e^{q_N/b}}$, and a buyer pays the
cost-function difference. Two properties made it the default
mechanism for thin markets: the maker's worst-case loss is bounded
at $b\ln 2$ (a fixed subsidy), and it \emph{always} quotes a price
--- no depth, no empty book. The trade-off: it is a house that
bleeds slowly to informed flow, which is exactly why live venues
migrated to order books.

\begin{example}[LMSR quotes slope up as you buy]\label{ex:lmsr}
\textbf{Strategy.} Instantiate a binary LMSR with $b=100$ and buy
two consecutive $50$-share YES blocks; compare the average fill to
an order book.

\textbf{Solution.} \texttt{lmsr\_demo.py} prints:
\begin{verbatim}
b=100: start price 0.500
buy 50 YES: paid 28.09  avg px 0.562  new YES px 0.622
buy 50 more: paid 33.92  avg px 0.678  new YES px 0.731
market-maker worst case = b*ln2 = 69.3
vs orderbook: LMSR always quotes (depth=b), book depth vanishes when you need it
\end{verbatim}
The first block fills at $0.562$ average, the second at $0.678$ ---
the AMM \emph{is} the book, and your own order walks you up its
price curve. The operator's exposure never exceeds
$b\ln 2 \approx 69.3$.

\textbf{Significance.} LMSR is the missing link between the two
books: it is an AMM (like the constant-product pools of the next
chapter) specialised to binary payout space, with a bounded-loss
guarantee no order book offers. When a venue must quote on 2000
long-tail questions at once, a $b\ln2$ subsidy per market is
cheaper than hiring a market-maker for each.
\end{example}

\section{Common mistakes}
\begin{itemize}
\item \textbf{Reading price as certainty.} $p=0.9$ is a probability, not a promise; one in ten
  90c favourites still loses, and tails cluster (fig.\ \ref{fig:pmhist}).
\item \textbf{Summing YES and NO sloppily.} A 3c sum gap on a binary market is usually vig or
  stale quotes, not an arb---check depth and fees before declaring free money.
\item \textbf{Scoring with accuracy.} Calling a 0.6 event ``yes'' and scoring 1/0 rewards
  overconfidence; use a proper rule to compare probability emitters.
\item \textbf{Trusting ECE alone.} Calibration without sharpness is cheap; a forecaster can
  be calibrated by always saying 0.5.
\item \textbf{Ignoring the resolution spec.} ``Xi out before 2027'' resolves on a specific
  oracle event---read the market's rules, not the headline.
\item \textbf{Forgetting frictions at the tails.} A 1c misprice is a 100x edge on paper and
  usually untradeable at size.
\item \textbf{Arbing the mid.} YES$+$NO sums to \$1 by
construction at the mid. At the \emph{ask} it costs a median
\$1.03 on live books (Example~\ref{ex:lock}) --- an arb that only
exists on the wrong side of the quote.
\item \textbf{Assuming every market quotes.} 60 of 100 snapshot
markets have no live two-sided book at all. ``The price'' of an
unquoted market is a stale print, not a level you can trade.
\end{itemize}

\section{Exercises}
\begin{exercise}[easy]
A YES share costs $0.35$. What does the NO share cost, and what does the sum imply about
the venue's vig?
\end{exercise}
\begin{solution}
$1-0.35=0.65$ in a frictionless book; a sum above \$1 is the effective vig on the pair.
\end{solution}

\begin{exercise}[easy]
You buy 100 YES shares at $0.62$ and the event resolves YES. PnL? At $q=0.7$, EV?
\end{exercise}
\begin{solution}
PnL $=100(1-0.62)=+38$; EV $=100(0.7-0.62)=+8$.
\end{solution}

\begin{exercise}[easy]
Brier score of forecaster who said $p=0.9$ on an event that failed? Log score (nats)?
\end{exercise}
\begin{solution}
Brier $=(0.9-0)^2=0.81$; log $=\log(1-0.9)=\log 0.1=-2.303$.
\end{solution}

\begin{exercise}[easy]
Why does the histogram in fig.\ \ref{fig:pmhist} have a U-shape?
\end{exercise}
\begin{solution}
Most-traded markets concentrate near resolution or longshots; mid-probability contests are
rarer among the top-volume snapshot.
\end{solution}

\begin{exercise}[easy]
A market's YES trades $0.52$ while its NO trades $0.44$. Sum?
\end{exercise}
\begin{solution}
$0.96$: buy both for $0.96$, guaranteed $1.00$ payout---a 4c theoretical arb, modulo
fees/resolution risk.
\end{solution}

\begin{exercise}
Show the log score is strictly proper: for truth $q$, expected score
$q\log p + (1-q)\log(1-p)$ is maximised at $p=q$.
\end{exercise}
\begin{solution}
Derivative $q/p-(1-q)/(1-p)=0\Rightarrow p=q$; second derivative negative---unique max at
truth.
\end{solution}

\begin{exercise}
From \texttt{data/pm\_markets.json}: report the 5 markets closest to $p=0.5$ and argue which
is the most contested.
\end{exercise}
\begin{solution}
Sort by $|p-0.5|$; the nearest-to-0.5 with high volume is the most contested---your script's
output is the answer.
\end{solution}

\begin{exercise}
A forecaster issues $p$ uniform on $[0.4,0.6]$ while true $q\in\{0.3,0.7\}$ equally. Show it is
uncalibrated even though mean forecast $=$ mean outcome.
\end{exercise}
\begin{solution}
Within any bin around 0.5, empirical rate $\neq\bar f$: forecasts near 0.6 still cover both
$q$'s, so conditional frequencies sit off the diagonal---matched marginals, mismatched
conditionals.
\end{solution}

\begin{exercise}
Why do sharp-tailed binary markets resist the ``$p$ is a probability'' reading more than
mid-range ones?
\end{exercise}
\begin{solution}
At $p<0.05$ the friction terms (fees, latency, thin book, resolution ambiguity) are the same
size as the price; at $p=0.5$ they are noise.
\end{solution}

\begin{exercise}
Dutch-book a three-outcome market: $p_A{=}0.4$, $p_B{=}0.35$, $p_C{=}0.30$. Free profit per
set?
\end{exercise}
\begin{solution}
Sum $=1.05$: sell all three, collect $1.05$, pay $1.00$ at resolution $=+0.05$ per set
before fees.
\end{solution}

\begin{exercise}
Give the reliability-diagram shape of (i) an overconfident forecaster, (ii) an underconfident
one, (iii) a well-calibrated one.
\end{exercise}
\begin{solution}
(i) S-curve bowed toward the middle (this chapter's fig.); (ii) inverted bow away from the
middle; (iii) on the diagonal within binning noise.
\end{solution}

\begin{exercise}
Explain why the log score's $-\infty$ penalty on a false 0\% is a \emph{feature}, not a bug.
\end{exercise}
\begin{solution}
It rules out certainty rhetoric: any forecaster who claims impossibility and is wrong even
once is eliminated---properness requires credence $>0$ on anything not logically excluded.
\end{solution}

\begin{exercise}[hard]
Derive the buy-YES Kelly fraction for edge $e=q-p$ at price $p$; plug in $p{=}0.62$, $q{=}0.70$.
\end{exercise}
\begin{solution}
Odds: win $(1-p)/p$ per unit; Kelly $f^*=(q(1-p)-p(1-q))/((1-p)p)\cdot p=(q-p)/(1-p)
=0.08/0.38\approx0.21$ of bankroll.
\end{solution}

\begin{exercise}[hard]
Extend \texttt{calibration.py}: estimate ECE with an equal-\emph{count} binning instead of
equal-width; explain when each lies.
\end{exercise}
\begin{solution}
Equal-count bins overweight sparse tails with noisy $\bar y$; equal-width bins let small bins
dominate the tails. Report both---your numbers are the answer.
\end{solution}

\begin{exercise}[hard]
Two forecasters: A issues $(0.9,0.8,0.7)$ on $(1,1,0)$; B issues $(0.7,0.7,0.7)$ on the same.
Score both under Brier and log; who is preferable and why does the answer differ?
\end{exercise}
\begin{solution}
A: Brier $(0.01+0.04+0.49)/3=0.18$, log $(\log.9+\log.8+\log.3)/3\approx-0.527$;
B: Brier $(0.09+0.09+0.49)/3=0.223$, log $3\log.7/3=-0.357$. A wins Brier (better resolution)
but loses log (the confident wrong 0.7 costs dearly)---the log rule punishes false confidence
harsher.
\end{solution}

\begin{exercise}[easy]
The cheapest YES$+$NO bundle costs $\$1.001$ at the ask. Compute
the guaranteed P\&L per \$1000 deployed.
\end{exercise}
\begin{solution}
$-\$1$: you pay $\$1001$ to receive a certain $\$1000$ --- a
riskless $-0.1\%$.
\end{solution}
\begin{exercise}[medium]
Why is the mid-price YES$+$NO sum exactly 1 on every market, yet
the ask-side bundle is not? Describe in terms of who wrote the
two quotes.
\end{exercise}
\begin{solution}
Outcome prices are complementary by construction (YES$=1-$NO at
mid). At the ask you are lifting \emph{two different makers'}
quotes --- the YES-ask sits above YES-mid and the NO-ask above
NO-mid, so the sum exceeds 1 by half of each spread.
\end{solution}
\begin{exercise}[hard]
You find a market where YES-ask $=0.48$ and NO-ask $=0.50$ ---
a lock costing $0.98$. List three ways this trade still loses or
fails, ordered by likelihood.
\end{exercise}
\begin{solution}
(a) The quotes are stale/unfillable at size (the 60\%-unquoted
majority); (b) resolution risk --- UMA dispute or ambiguous
wording can price YES$+$NO off \$1; (c) settlement delay ties up
collateral against a $2\%$ certain gain that may take months.
\end{solution}

\section*{Notes}
Prediction markets are the cleanest place to learn probabilistic thinking: binary payoffs,
public prices, and resolution oracles that grade every forecast for you. The same machinery
(proper scores, calibration, sum-to-one consistency) reappears when evaluating ML classifiers
and trading signals---a mispriced market and a miscalibrated model fail the same way.

\section*{Further reading}
\begin{itemize}
\item Polymarket documentation --- market mechanics, resolution and the API used in this
chapter:

\url{https://docs.polymarket.com/}
\item Wikipedia, ``Prediction market'' --- history, mechanism design, and limits:

\url{https://en.wikipedia.org/wiki/Prediction_market}
\item Wikipedia, ``Scoring rule'' --- the propriety theory behind \S\ref{sec:scoring}:

\url{https://en.wikipedia.org/wiki/Scoring_rule}
\item Wikipedia, ``Brier score'' --- decomposition (reliability/resolution/uncertainty):

\url{https://en.wikipedia.org/wiki/Brier_score}
\item FiveThirtyEight, ``Checking Our Work'' --- calibration analysis on real public
forecasts:

\url{https://projects.fivethirtyeight.com/checking-our-work/}
\item Polymarket Gamma API --- the public endpoint used by \texttt{arb\_check.py}:

\url{https://gamma-api.polymarket.com/}
\end{itemize}

\begin{rulebox}[Watch instead]
\begin{itemize}
\item ``What Is POLYMARKET? | The Crypto Prediction Market Easy Explained'' --- What The
Asset (9:43). Contract mechanics and the venue used in this chapter.
\url{https://www.youtube.com/watch?v=Vqwjwj0-nmQ} (accessed 2026-10-11)
\item ``How Prediction Markets Work (and why you should care)'' --- Catalyst (7:05). Price-as-
probability reading; matches \S\ref{sec:pmcontract}.
\url{https://www.youtube.com/watch?v=tWFuBuo_B10} (accessed 2026-10-11)
\end{itemize}
\end{rulebox}
